\documentclass[11pt]{article}
\usepackage[T1]{fontenc}
\usepackage[margin=1in]{geometry}
\usepackage{amsmath,amssymb,amsthm,mathtools,hyperref}
\hypersetup{hidelinks}
\newtheorem{theorem}{Theorem}[section]
\newtheorem{lemma}[theorem]{Lemma}
\newtheorem{proposition}[theorem]{Proposition}
\newtheorem{corollary}[theorem]{Corollary}
\theoremstyle{definition}
\newtheorem{definition}[theorem]{Definition}
\newtheorem{remark}[theorem]{Remark}
\newcommand{\Q}{\mathbf Q}
\newcommand{\C}{\mathbf C}
\newcommand{\F}{\mathbf F}
\newcommand{\Z}{\mathbf Z}
\newcommand{\CH}{\operatorname{CH}}
\newcommand{\Pic}{\operatorname{Pic}}
\newcommand{\Alt}{\operatorname{Alt}}
\newcommand{\cl}{\operatorname{cl}}
\newcommand{\id}{\operatorname{id}}
\newcommand{\Spec}{\operatorname{Spec}}
\newcommand{\im}{\operatorname{im}}
\title{Balanced mixed determinants of prime cyclic covers:\\a revised fusion argument}
\author{Kyle Kistner}
\date{4 September 2026}
\begin{document}
\maketitle
\begin{abstract}
We isolate a proposed consequence of the prime balanced-tuple theorem and
Schoen's simple-tuple cycle theorem. For an arbitrary finite list of positive-genus
cyclic covers of the projective line of the same odd prime degree, balance of the
combined branch residues implies algebraicity of a specified rational space of
mixed determinant classes. The argument constructs a compact-type degeneration,
uses algebraic smoothing and specialization, and treats disconnected pairing
graphs by an explicit rational correspondence. This is a classical proof draft
using stated external results, not a machine-checked theorem. No assertion about
the full Hodge ring, a Mumford--Tate group, or novelty follows without additional
work; the all-powers application is stated separately with its missing input.
\end{abstract}

\section{Scope and conventions}
Fix an odd prime $p$, a primitive root $\zeta=\exp(2\pi i/p)$, and
$K=\Q(\zeta)$. Let $N\geq1$. For each $1\leq j\leq N$, let
\[
 f_j:Y_j\longrightarrow\mathbf P^1_\C
\]
be a smooth connected cyclic cover of degree $p$, with a specified faithful
generator $\sigma_j$ of its deck group. Its ordered nonzero branch residues are
$a_{j1},\ldots,a_{js_j}\in\F_p^\times$, where $s_j\geq3$ and
$\sum_i a_{ji}=0$ in $\F_p$. The residue $a$ means that positive local
monodromy is $\sigma_j^a$. Set
\[
 A_j=J(Y_j),\quad V_j=H^1(A_j,\Q)=H^1(Y_j,\Q),\quad
 h_j=s_j-2,\quad S=\sum_j s_j,\quad H=\sum_j h_j=S-2N.
\]
We use the induced action on cohomology for the chosen $K$-action, sending
$\zeta$ to $\sigma_j^*$. For example the corresponding automorphism of
$\Pic^0(Y_j)$ can be taken to be pullback by $\sigma_j^{-1}$, so that the
Abel--Jacobi identification on $H^1$ has exactly this convention.

The genus is $g_j=(p-1)h_j/2$. Moreover $V_j$ is a $K$-vector space of
rank $h_j$: its deck-invariant part is $H^1(\mathbf P^1,\Q)=0$, and
$\Q[C_p]=\Q\times K$. All $h_j$ are positive. Repetitions among the curves
and different choices of generators on the same curve are permitted.

Write $[a]_p\in\{1,\ldots,p-1\}$ for the representative of a nonzero residue.
The balance condition is
\begin{equation}\label{eq:balance}
 2\sum_{j,i}[b a_{ji}]_p=pS\qquad\text{for every }b\in\F_p^\times.
\end{equation}
It implies that $S$, and hence $H$, is even. Since $H>0$, $H\geq2$.

For an embedding $\tau_b:K\hookrightarrow\C$ with $\tau_b(\zeta)=\zeta^b$,
write $V_{j,b}$ for the associated eigenspace. There is a canonical rational
subspace $W\subset H^H(\prod_j A_j,\Q)$ specified after extension of scalars by
\begin{equation}\label{eq:W}
 W_\C=\bigoplus_{b\in\F_p^\times}
       \bigotimes_{j=1}^N\bigwedge^{h_j}V_{j,b}.
\end{equation}
The embedding is the usual exterior-algebra and K\"unneth embedding. Each
summand has dimension one, and $\dim_\Q W=p-1$. Abstractly,
\[
 W\simeq\bigwedge_K^H\Bigl(\bigoplus_jV_j\Bigr)
       \simeq\bigotimes_{K,j}\bigwedge_K^{h_j}V_j.
\]
These formulas mean the rational realization \eqref{eq:W}; they do not mean the
full tensor product over $\Q$ of the individual determinant spaces.

\begin{theorem}[Standalone balanced fusion statement]\label{thm:fusion}
With the hypotheses above, assume \eqref{eq:balance}. Then
\[
 W\subset\im\left(\cl:\CH^{H/2}\left(\prod_j A_j\right)_\Q
                       \longrightarrow H^H\left(\prod_j A_j,\Q\right)\right).
\]
In particular $W(H/2)$ is a rational Hodge structure of type $(0,0)$ generated
by classes of algebraic cycles. No genericity of the branch points, full
monodromy, or endomorphism-algebra hypothesis is required.
\end{theorem}

\paragraph{Status.}
The following proof reduces this statement to the precise published inputs
in Section~\ref{sec:inputs}. The project-specific rational projector, residue
conversion, gluing, and reduction to those inputs are supplied below. This draft
has not received independent mathematical referee review; finite checks and the
repository's Lean layer do not certify its algebraic geometry.

\section{Rational realization and algebraic projectors}\label{sec:projectors}
We first make the tensor product over $K$ usable on algebraic cycles. Let
$U_1,\ldots,U_m$ be copies of the $V_j$, with their specified deck operators
$T_r$. On $\bigotimes_{\Q,r}U_r$, set
\begin{equation}\label{eq:projector}
 e_{uv}=\frac1p\sum_{a=0}^{p-1}T_u^aT_v^{-a},\qquad
 E_m=\prod_{r=2}^m e_{1r},\qquad E_1=\id.
\end{equation}
All operators in these formulas act on their indicated tensor slots.

\begin{lemma}\label{lem:slots}
The rational operator $E_m$ is an idempotent with image
$\bigoplus_b\bigotimes_r U_{r,b}$. It is induced by an algebraic
correspondence on the product containing the tensor slots. For $m$ identical
slots, it commutes with every ordinary tensor permutation.
\end{lemma}
\begin{proof}
Over a splitting field, the value of $e_{uv}$ on an eigentensor with labels
$b_u,b_v$ is
\[
 \frac1p\sum_{a=0}^{p-1}\zeta^{a(b_u-b_v)}
 =\begin{cases}1&b_u=b_v,\\0&b_u\ne b_v.\end{cases}
\]
This proves the image and idempotence. Equality of all labels is invariant under
permutations, proving the commutation assertion after scalar extension and
hence over $\Q$. Each $T_r^a$ comes from a deck automorphism on an abelian
factor. Their graphs, rational linear combinations, and compositions are
algebraic correspondences. The $H^1$ K\"unneth projection is also algebraic:
for an abelian variety of dimension $g$ and an integer $n\geq2$, the polynomial
\[
 \pi_k=\prod_{0\leq d\leq 2g,\ d\ne k}
               \frac{[n]^*-n^d}{n^k-n^d}
\]
acts as the projection to $H^k$. It is represented cohomologically by a
rational algebraic correspondence. Only this cohomological assertion about
$\pi_k$ is needed, not an equality of Chow motives.
\end{proof}

For an abelian variety $A$ and $m\geq1$, let $\mu:A^m\to A$ be addition,
$\delta:A\to A^m$ the diagonal, and
$P=\pi_1^{\otimes m}$. The map
\[
 \iota_m=P\mu^*:H^m(A,\Q)\longrightarrow H^1(A,\Q)^{\otimes m}
\]
is the unnormalized alternating inclusion. The cup product
$q_m=\delta^*|_{H^1(A)^{\otimes m}}$ satisfies
\begin{equation}\label{eq:wedge}
 q_m\iota_m=m!\,\id,\qquad
 \iota_mq_m=m!\,\Alt_m,
 \quad \Alt_m=\frac1{m!}\sum_{\pi\in\mathfrak S_m}
                  \operatorname{sgn}(\pi)P_\pi.
\end{equation}
Indeed expansion of $\mu^*v=\sum_r\operatorname{pr}_r^*v$ for $v\in H^1(A)$
gives the first identity; the exterior-algebra relation gives the second.
Geometric permutations on products of odd cohomology classes have Koszul
signs. Thus $P_\pi$ here is the ordinary tensor permutation, realized by the
geometric permutation correspondence multiplied by $\operatorname{sgn}(\pi)$.
There is no hidden assumption that these two actions coincide.

For $A=A_j$ and $m=h_j$, the correspondence
\begin{equation}\label{eq:Dprojector}
 Q_j=\frac1{h_j!}q_{h_j}E_{h_j}\iota_{h_j}
\end{equation}
is an idempotent on $H^{h_j}(A_j)$ with image
$D_j=\bigoplus_b\bigwedge^{h_j}V_{j,b}$. Idempotence follows from
\eqref{eq:wedge} and Lemma~\ref{lem:slots}. This proves rationality and
algebraic realizability of the determinant summands without selecting a single
complex embedding as a rational subspace.

\begin{remark}[The scalar action must be induced from one slot]\label{rem:scalar}
On $D_j$, the scalar action of $k\in K$ is the action obtained from $k$ on
\emph{one} slot after applying $E_{h_j}$ and alternation. All slot actions then
agree, so this defines a unital $\Q$-algebra action. The corresponding scalar
$\zeta$ acts by $\zeta^b$ on the $b$ determinant line. In contrast,
$\bigwedge^{h_j}(\sigma_j^*)$ acts by $\zeta^{b h_j}$. These are different
operators, and the latter is trivial when $p\mid h_j$. One must not use this
exterior-power deck action to identify embeddings between determinant factors.
Our construction needs no hypothesis $\gcd(p,h_j)=1$.
\end{remark}

Apply $E_H$ to all raw $H^1$ slots simultaneously, alternate within the $j$-th
block of $h_j$ slots, and use \eqref{eq:wedge} to return to
$H^H(\prod_j A_j)$. This gives a rational algebraic projector $Q_W$ with image
$W$. More explicitly, put $\iota=\bigotimes_j\iota_{h_j}$,
$q=\bigotimes_j q_{h_j}$, and $c=\prod_jh_j!$. On the indicated
K\"unneth degree $(h_1,\ldots,h_N)$ the formula is
\begin{equation}\label{eq:globalprojector}
 Q_W=c^{-1}qE_H\iota;
\end{equation}
extend it by zero on the other degrees using algebraic K\"unneth projectors.
The same calculation proves idempotence. Its action takes algebraic classes
to algebraic classes and fixes $W$.

\section{Exact external inputs}\label{sec:inputs}
The proof uses the following results; none is asserted as a new theorem here.
\begin{enumerate}
\item \textbf{Prime balanced tuples.} A tuple of nonzero residues of even
length at least four, with sum zero modulo the odd prime $p$, satisfying
\eqref{eq:balance}, can be partitioned into pairs $(a,-a)$. This is the
prime equality $B_p^n=D_p^n$ stated on p.~24 before Theorem A of
Aoki~\cite{Aoki}, where the prime case is credited to W.~Parry; here
$n=S-2$ is positive and even. The 1984 erratum corrects the introduction's
Theorem B, not this input.
\item \textbf{Simple-tuple cycles.} For a smooth connected cyclic curve with
faithful action and simple branch tuple (a union of opposite pairs), Schoen's
Theorem 2.0 and Corollary 3.1~\cite{Schoen} give algebraic classes spanning the
primitive determinant space on its primitive Jacobian factor. We use this in
characteristic zero, where Theorem 2.0 also holds over an arbitrary algebraically
closed field. For a prime cover of $\mathbf P^1$ the primitive factor is the
whole Jacobian up to isogeny, as shown in Section 1. The transfer in Corollary
3.1 uses the algebraicity of inverse Lefschetz on abelian varieties; its argument
works over algebraically closed characteristic-zero fields as well, by descent
and comparison. No claim about arbitrary abelian varieties of Weil type is used.
\item \textbf{Balanced admissible-cover deformation.} Abramovich--Corti--Vistoli
\cite{ACV}, Theorems 3.0.2 and 4.3.2 and Corollary 3.0.5, identify balanced
admissible $G$-covers with balanced twisted covers and give their unobstructed
deformation theory, with a flat map to the moduli stack of stable pointed
target curves. Section 3 identifies the balanced node smoothing directions.
\item \textbf{Compact-type Picard and specialization.} For a proper nodal
curve family whose fibers have tree dual graph, $\Pic^0$ is an abelian scheme,
with special fiber the product of the component Jacobians. The relative Picard
input is the semistable Picard theorem attributed to Deligne in
\cite[\S9.4, Theorem 1, p.~259]{BLR}; the special-fiber description is
\cite[\S9.2, Example 8, pp.~246--247]{BLR}. An explicitly inspected primary use of precisely
this relative assertion is Achter--Pries~\cite[\S3.1]{AP}.
For a smooth proper scheme over a DVR, specialization of Chow classes is
compatible with smooth-proper specialization in $\ell$-adic cohomology.
For the cycle map see SGA~6, X.7.13, especially 7.13.11, and X.7.16,
\cite[pp.~576--579, 585]{SGA6}. Section 7.13 assumes a henselian trait;
we make this base change explicitly in Section~\ref{sec:specialize}.
Schoen~\cite[pp.~30--31]{Schoen} invokes the same compatibility.
The cycle-level construction and base change
are recorded in \cite[Tags 0H4H, 0H4M]{Stacks}.
\end{enumerate}

\paragraph{Source-audit boundary.}
Aoki's exact prime assertion and erratum were pinned by the earlier repository
audit. In this repair the indicated Schoen statements, his specialization
paragraph, the 1998 addendum, ACV's deformation argument and equivalence, the
Achter--Pries Picard assertion, and the Stacks cycle construction were inspected.
The final cross-review directly inspected the original BLR pages 246--247 and
258--259, and the original SGA~6 pages 576--579 and 585. The latter verify
the henselian assumption, smooth-proper comparison, rational cycle-class
compatibility in 7.13.11, and passage to geometric fibers in 7.16. We do not
claim to have inspected the scan's missing sections 7.14--7.15. These are
standard imported results with the precise forms used here stated above,
not new theorems proved by this manuscript.

\section{From balance to a compact-type cover}\label{sec:glue}
Choose a pairing of all $S$ branch occurrences into opposite residues using
the first input. Form a multigraph $\Gamma$ with vertices $1,\ldots,N$ and
one edge for every chosen pair; both ends can lie on the same vertex. An
occurrence is used once. Let $I$ be a connected component of this graph.
If $|I|>1$, choose a spanning tree $T_I$ from its non-loop edges; if $|I|=1$,
use the empty tree. Every branch occurrence on a vertex of $I$ is paired with
another occurrence in $I$.

\begin{lemma}[Residue versus tangent character]\label{lem:tangent}
At a totally ramified branch point with residue $a$, the action of $\sigma_j$
on a local tangent coordinate is multiplication by $\zeta^{a^{-1}}$.
Consequently residues $a$ and $-a$ have inverse tangent characters for the
chosen common deck generator.
\end{lemma}
\begin{proof}
Write the cover locally as $x=u^p$. Positive monodromy about $x=0$ sends
$u$ to $\zeta u$. By definition this is $\sigma_j^a$, so if
$\sigma_j(u)=\zeta^c u$ then $ac=1$ modulo $p$. Also
$(-a)^{-1}=-a^{-1}$.
\end{proof}

\begin{lemma}[Balanced gluing]\label{lem:glue}
For every $I$ there is a connected balanced admissible $C_p$-cover
$Y_{I,0}\to B_{I,0}$ whose component covers are exactly the chosen
$Y_j\to\mathbf P^1$ for $j\in I$, whose source and target dual graphs are
$T_I$, and whose remaining branch tuple is simple. Its source has compact
type and genus $\sum_{j\in I}g_j$.
\end{lemma}
\begin{proof}
For each tree edge glue its two indicated base branch points, and glue their
unique ramification points upstairs. The deck action agrees under gluing,
since the glued points are fixed. The curve maps glue to a finite morphism
of nodal curves; away from markings and nodes they are principal $C_p$-bundles.
Lemma~\ref{lem:tangent} gives the balanced action at every node. Both graphs
are the chosen tree because a totally ramified point has a unique preimage.

Each target component has exactly $s_j$ special points: replacing a selected
branch marking by a node changes none of their number. Hence the pointed
target is stable. Each source component has positive genus. A genus-one
component in a non-singleton tree has a node and is stable; a singleton graph
component has even $s_j\geq4$, hence $g_j\geq2$. Thus the source is stable
after forgetting branch markings as well. The cover satisfies the admissible
$G$-cover definition of \cite[Definition 4.3.1]{ACV}.

Only complete pairs belonging to $T_I$ were removed. The unremoved tuple is
therefore a union of opposite pairs and has length
\begin{equation}\label{eq:counts}
 s_I=\sum_{j\in I}s_j-2(|I|-1)=H_I+2,\qquad
 H_I=\sum_{j\in I}h_j.
\end{equation}
All original occurrences within $I$ are paired, so their total number is even;
thus $H_I$ is positive and even, $H_I\geq2$, and $s_I\geq4$. The tree graph
has first Betti number zero, giving both compact type and the asserted genus.
\end{proof}

\begin{lemma}[Algebraic simultaneous smoothing]\label{lem:smooth}
The cover of Lemma~\ref{lem:glue} is the closed fiber of an algebraic family
of admissible $C_p$-covers over the spectrum of a DVR essentially of finite
type over $\C$. Its geometric generic fiber is a smooth connected cyclic
cover of $\mathbf P^1$ with the simple remaining branch tuple. After a finite
base change, the family has an abelian scheme
$\mathcal A_I=\Pic^0(\mathcal Y_I/R)$, equipped with the chosen $K$-action,
whose special fiber is $\prod_{j\in I}A_j$.
\end{lemma}
\begin{proof}
At a glued node of residue $a$ the local smoothing is
\[
 uv=t,\quad x=u^p,\quad y=v^p,\quad xy=t^p,\qquad
 \sigma(u,v)=(\zeta^{a^{-1}}u,\zeta^{-a^{-1}}v).
\]
This is equivariant and smooths the node when $t\ne0$. Notice the power
$t^p$ in the target parameter. The cover map is finite; it is not asserted
flat at a node. Both nodal curve families are flat over the parameter base.

By the ACV input these local balanced deformations globalize in an algebraic
stack with smooth local charts. More precisely the local-to-global deformation
sequence for a nodal twisted curve surjects onto its local node deformation
spaces: the obstruction to this surjection is $H^2$ of its tangent sheaf,
which is zero for a curve; the node deformation sheaf is supported at finitely
many points. The invariant smoothing line survives exactly because the tangent
characters are inverse. The torsor adds no obstruction, as its target $BG$ has
zero cotangent complex in characteristic zero. These are the deformation
calculations in ACV Section 3. Thus no node parameter is forced to vanish.

Take an algebraic smooth atlas at our point and a curve through it with tangent
direction outside the finitely many node-parameter hyperplanes. Such a curve
exists: choose a suitable line in local \'etale coordinates and take its
inverse image. Normalize and localize at the point, obtaining a DVR. Each
source node parameter has a simple zero, so the source total space is regular
at these nodes. All node parameters are nonzero at the generic point.
Shrinking the atlas if
necessary, smoothness away from the specified nodes is open, so the generic
source and target are smooth. This constructs an algebraic family with the
prescribed closed fiber, not merely a formal or analytic plumbing. The empty
tree case can instead use a constant family over a DVR.

The source is geometrically connected: its connected special fiber has
$h^0(\mathcal O)=1$, which bounds the generic value from above by
semicontinuity, while constants bound it from below. The target has genus
zero and the remaining marked sections; after a finite extension if needed
it is $\mathbf P^1$. The labeled inertia data at these sections is locally
constant, so the generic tuple is exactly that in \eqref{eq:counts}.

There are only generic and closed geometric fibers over the DVR; both have
compact type. The compact-type Picard input gives the abelian scheme. For
completeness, on the closed fiber normalization gives a Picard extension with
torus of rank $b_1(T_I)=0$. Restriction of line bundles therefore identifies
its identity component with the product of the component Jacobians, compatibly
with deck actions. Since $1+\sigma+\cdots+\sigma^{p-1}$ acts as zero on the
generic Jacobian, it acts as zero on the abelian scheme: a homomorphism of
separated group schemes equal to zero on the dense generic fiber is zero.
Thus the deck action factors through $K$ in rational endomorphisms everywhere.
\end{proof}

\section{Spreading cycles and specializing their classes}\label{sec:specialize}
Fix $I$, let $L$ be the fraction field of the DVR in
Lemma~\ref{lem:smooth}, and fix an algebraic closure $\overline L$. The
geometric generic curve has $K$-rank $H_I$ in $H^1$, by
\eqref{eq:counts}. The simple-tuple theorem gives cycle classes spanning
\[
 W_{I,\overline\eta}
 =\bigwedge_{K\otimes_\Q\Q_\ell}^{H_I}
       H^1(\mathcal A_{I,\overline\eta},\Q_\ell)
 \subset H^{H_I}(\mathcal A_{I,\overline\eta},\Q_\ell).
\]
Here and below the wedge notation denotes the matching-embedding realization
after extension to $\Q_\ell$, including the appropriate Tate twist when
comparing cycle classes. Fix any rational prime $\ell$.

\begin{lemma}\label{lem:specialization}
The space
\[
 W_I=\bigoplus_b\bigotimes_{j\in I}\bigwedge^{h_j}V_{j,b}
 \subset H^{H_I}\left(\prod_{j\in I} A_j,\Q\right)
\]
is generated by algebraic classes.
\end{lemma}
\begin{proof}
Choose finitely many rational algebraic cycles whose $\Q_\ell$-classes span the
geometric generic determinant space, applying the generic-fiber algebraic
determinant projector first if the supplied cycles have other components.
Each cycle is defined over a finite
extension of $L$: its finitely many defining equations have coefficients
algebraic over $L$. A single finite extension $L'/L$ suffices for all chosen
cycles. The DVR is excellent, so its integral closure in $L'$ is finite;
localize at a point above the closed point to obtain a DVR $R'$. Its residue
field is $\C$. Pulling back the abelian scheme leaves its complex special
fiber unchanged.

Take the closures of the generic cycles in this proper smooth scheme and
specialize in the Chow group. An integral closure of a generic cycle
dominates the DVR and is flat over it; its special fiber with multiplicities
gives the specialization. For the cohomological comparison, now pass from
$R'$ to its strict henselization $R'^{\mathrm{sh}}$, with a compatible choice
of geometric generic point. Its residue field is still $\C$, so its closed
fiber and the specialized complex cycles are unchanged. Pull back the abelian
scheme and the cycles to this trait. Specialization commutes with this flat
DVR base change; all closures and finite-extension descent were already
constructed over $R'$. The henselian hypothesis in SGA~6, X.7.13 is now met.
The compatibility cited in Section~\ref{sec:inputs}
identifies these classes with the original generic classes under the
smooth-proper $\ell$-adic specialization isomorphism. This is a specialization
of cycles defined over a finite extension, not an assumption that an arbitrary
Hodge class is algebraic because it lies in a family.

The relative $H^1$ local system carries the deck action. The projectors from
Section~\ref{sec:projectors}, formed from addition, diagonals, multiplication,
and deck maps of the abelian scheme, commute with specialization. Thus
specialization identifies the generic determinant space with
$\bigwedge_K^{H_I}(\bigoplus_{j\in I}V_j)$ on the closed fiber. Expanding the
top exterior power of each $b$-eigenspace gives exactly $W_I$: to reach degree
$H_I=\sum h_j$ one must use all $h_j$ dimensions of each summand. In
particular no other K\"unneth degrees occur.

Apply the algebraic special-fiber projector if necessary so that every
specialized cycle class lies in $W_I$. Their rational Betti classes have
$\Q_\ell$-span $W_I\otimes\Q_\ell$. Betti--\'etale comparison and extension
of scalars then show that their $\Q$-span is $W_I$: the quotient of $W_I$ by
that span becomes zero over the field $\Q_\ell$, and is therefore zero.
This proves rational, not merely $\ell$-adic, algebraicity.
\end{proof}

\begin{proof}[Proof of Theorem~\ref{thm:fusion}]
For every connected component $I$ of the pairing graph,
Lemma~\ref{lem:specialization} supplies algebraic classes spanning $W_I$.
Take their external products. They span $\bigotimes_{\Q,I}W_I$, which
contains independent embedding choices on the different graph components.
The projector $Q_W$ of \eqref{eq:globalprojector} kills precisely the terms
whose embedding choices disagree. Its image on that tensor product is $W$,
since every matching term is fixed. Applying this algebraic correspondence
to the external-product cycles proves the theorem. If the graph is connected,
this last projection is redundant; if it is disconnected, it is essential.
\end{proof}

\section{Independent Hodge-type check and example}
The Hodge type can also be checked without a monodromy theorem. In a coordinate
with infinity unramified write $y^p=\prod_i(x-t_i)^{a_i}$ with
$1\leq a_i\leq p-1$ and $\sum a_i=kp$. For $l=p-b$, forms of deck eigenvalue
$\zeta^b$ have the shape $r(x)dx/y^l$. Holomorphy at $t_i$ forces
$r(x)$ divisible by $(x-t_i)^{\lfloor la_i/p\rfloor}$, since the valuation is
$p\operatorname{ord}_{t_i}(r)+p-1-la_i$. Holomorphy at infinity bounds the
remaining polynomial degree by $lk-\sum_i\lfloor la_i/p\rfloor-2$.
Consequently
\[
 \dim V_{j,b}^{1,0}
 =-1+\frac1p\sum_i[-b a_{ji}]_p
 =s_j-1-\frac1p\sum_i[b a_{ji}]_p.
\]
The fractional-part sum is a positive integer, so the formula includes the
zero-dimensional case. Under balance, the sum over $j$ is
$S-N-S/2=H/2$. Each line of \eqref{eq:W} therefore has type $(H/2,H/2)$.
This also fixes the sign convention rather than importing an ambiguous
Chevalley--Weil formula.

For $p=5$, take branch tuples
\[
 (1,1,3),\qquad (1,2,4,4,4).
\]
They give $K$-ranks $1,3$ and genera $2,6$. Their $H^{1,0}$ multiplicities
for $b=1,2,3,4$ are respectively $(1,1,0,0)$ and $(1,1,2,2)$, whose sum is
$(2,2,2,2)$. The concatenation pairs into three $(1,4)$ pairs and one
$(3,2)$ pair. Glue the residue-$3$ point of the first curve to the residue-$2$
point of the second. The remaining tuple is $(1,1,1,4,4,4)$, the smooth
fused curve has genus $8$ and $K$-rank $4$, and the theorem gives a
four-dimensional rational space of codimension-two classes on $A_1\times A_2$.
Neither individual determinant space is of middle Hodge type. The mixed
balance is the relevant condition.

\section{The separate all-powers implication}\label{sec:conditional}
Let $A$ be an abelian variety assembled, up to isogeny, from cyclic quotient
Jacobians of odd prime degree. To deduce the Hodge conjecture for every power
of $A$, the following \emph{additional generator statement} suffices:

\medskip\noindent
\textbf{Generator hypothesis.} Every rational Hodge class on every $A^m$
is in the span obtained from divisor classes and spaces $W$ satisfying
\eqref{eq:balance} by external products, cup products, pullbacks and
pushforwards along algebraic correspondences, including the isogenies,
quotient maps, permutations, and rational projectors needed for realization.
The determinant words have positive multiplicities of the chosen cyclic
$H^1$ factors; conjugate words use inverse chosen generators. Any dualization
is implemented with its polarization and Tate twist before using this form.
\medskip

\begin{corollary}[Conditional all-powers application]
Assume the generator hypothesis for $A$. Then every rational Hodge class on
every power of $A$ is algebraic.
\end{corollary}
\begin{proof}
Divisor Hodge classes are algebraic by the Lefschetz $(1,1)$ theorem. Theorem
\ref{thm:fusion} handles each indicated determinant space. Every operation in
the generator hypothesis preserves algebraic cycle classes.
\end{proof}

The generator hypothesis is a distinct representation-theoretic and rational
descent problem. This manuscript does not prove the claimed full monodromy,
the treatment of compact embedding sectors, graph-arrow signs, the complete
central torus, or exhaustion by determinant words. An integer lattice basis
of determinant relations alone does not prove generation of the effective
tensor monoid; negative exponents need separately realized duals and twists,
or an appropriate monoid-generation argument. No unconditional all-powers
claim is made here.

\section{Relation to prior work and review gates}
The input for a single simple cyclic cover is Schoen's. Compact-type clutching
and algebraic specialization are also established tools; the proposed
contribution is their combination for an arbitrary list of separately
unbalanced prime cyclic factors whose combined signatures balance, including
the disconnected-graph projection. A bounded literature search did not locate
this exact quantified mixed-product formulation, but does not establish novelty.
It should be compared explicitly with the cyclic-cover clutching literature
and later work on exceptional classes on products of CM and moving factors.

Schoen's 1998 addendum~\cite{SchoenAddendum} corrects the missing polarization
invariant in the statement of the 1988 Theorem 3.2. It explicitly continues to
use Corollary 3.1. Our argument uses the latter simple-tuple input and does not
import the uncorrected general fourfold statement.

The next review gates are: independent expert verification of the geometric
proof beyond the internal cross-review; an exact prior-art comparison for the mixed-product
statement; and, for any global Hodge application, a proof of the separate
generator hypothesis. The finite script accompanying this revision checks
projector identities and tuple/graph examples only. It cannot verify
smoothability, cycle existence, or specialization.

\begin{thebibliography}{99}
\bibitem{Aoki} N. Aoki, \emph{On some arithmetic problems related to the Hodge
cycles on the Fermat varieties}, Math. Ann. 266 (1983), 23--54,
\url{https://doi.org/10.1007/BF01458703}. Erratum, Math. Ann. 267 (1984), 572,
\url{https://eudml.org/doc/163917}.
\bibitem{Schoen} C. Schoen, \emph{Hodge classes on self-products of a variety
with an automorphism}, Compositio Math. 65 (1988), 3--32,
\url{https://www.numdam.org/item/CM_1988__65_1_3_0.pdf}.
\bibitem{SchoenAddendum} C. Schoen, \emph{Addendum to: Hodge classes on
self-products of a variety with an automorphism}, Compositio Math. 114 (1998),
329--336, \url{https://doi.org/10.1023/A:1000566205021}.
\bibitem{ACV} D. Abramovich, A. Corti, and A. Vistoli, \emph{Twisted bundles
and admissible covers}, Comm. Algebra 31 (2003), 3547--3618; inspected version
arXiv:math/0106211v1, \url{https://arxiv.org/abs/math/0106211v1}.
\bibitem{BLR} S. Bosch, W. L\"utkebohmert, and M. Raynaud, \emph{N\'eron
Models}, Ergebnisse der Mathematik 21, Springer, 1990, Chapter 9,
\url{https://doi.org/10.1007/978-3-642-51438-8}.
\bibitem{AP} J. D. Achter and R. Pries, \emph{The integral monodromy of
hyperelliptic and trielliptic curves}, Math. Ann. 338 (2007), 187--206,
\url{https://www.math.colostate.edu/~achter/math/trigonal.pdf}.
\bibitem{SGA6} P. Berthelot, A. Grothendieck, and L. Illusie, eds.,
\emph{Th\'eorie des intersections et th\'eor\`eme de Riemann--Roch}, SGA 6,
Lecture Notes in Math. 225, Springer, 1971, Expos\'e X, 7.13 (especially
7.13.11) and 7.16, pp.~576--579 and 585,
\url{https://library.slmath.org/nonmsri/sga/sga/pdf/sga6.pdf}.
\bibitem{Stacks} The Stacks Project authors, \emph{Specialization of cycles},
Tags 0H4H and 0H4M, \url{https://stacks.math.columbia.edu/tag/0H4H} and
\url{https://stacks.math.columbia.edu/tag/0H4M}.
\end{thebibliography}
\end{document}
